\begin{figure}[H]
    \centering
    {
    \chemnameinit{\chemfig{^{-}OOC-[7]C(-[5]H)=C(-[7]{\color{red}H})-[1]{\color{blue}C}|{\color{blue}OO^{-}}}}
    \schemestart
        \chemname{\chemfig{^{-}OOC-[7]C(-[5]H)=C(-[7]{\color{red}H})-[1]{\color{blue}C}|{\color{blue}OO^{-}}}}{\small\bfseries Maleato}
        \arrow{<->>}
        \chemname{\chemfig{^{-}OOC-[7]C(-[5]H)=C(-[1]{\color{red}H})-[7]{\color{blue}C}|{\color{blue}OO^{-}}}}{\small\bfseries Fumarato}
    \schemestop
    \chemnameinit{}
    \vspace*{1em} \par}
    \rule{\linewidth}{0.2pt}
    \caption[Isomerización del maleato en fumarato por la maleato isomerasa]{Reacción catalizada por la maleato isomerasa (EC 5.2.1.1). Esta enzima lleva a cabo una reacción de isomeración, o sea, realiza un cambio en el ordenamiento de los átomos del sustrato, el maleato, y que genera un isómero del sustrato, fumarato.}
    \label{fig:reaccion_maleato_isomerasa}
\end{figure}